\Article{244}
Wenn durch eine Form $f$ die Zahl $a$, durch eine Form $f'$ die Zahl $a'$ dargestellt werden kann und die Form $F$ in $ff'$ transformierbar ist, so erkennt man ohne Mühe, dass das Produkt $aa'$ durch die Form $F$ darstellbar ist. Daraus folgt sofort: Sind die Determinanten dieser Formen negativ, so ist die Form $F$ positiv, wenn $f$, $f'$ entweder beide positiv oder beide negativ sind; dagegen wird $F$ negativ, wenn eine der Formen $f$, $f'$ positiv und die andere negativ ist. Betrachten wir besonders den Fall des vorigen Artikels, in dem $F$ aus $f$, $f'$ zusammengesetzt ist. Außerdem haben $f$, $f'$ und $F$ dieselbe Determinante $D$. Wir nehmen außerdem an, dass die Zahlen $a$, $a'$ durch die Formen $f$, $f'$ mit teilerfremden Werten der Unbestimmten dargestellt werden und die erste Darstellung zum Wert $b$ des Ausdrucks $\surd D\ (\mathrm{mod.}\ a)$, die zweite zum Wert $b'$ des Ausdrucks $\surd D\ (\mathrm{mod.}\ a')$ gehört, und setzen $bb-D=ac$, $b'b'-D=a'c'$. Dann sind nach Artikel 168 die Formen $(a,b,c)$, $(a',b',c')$ den Formen $f$, $f'$ eigentlich äquivalent; somit ist $F$ auch aus diesen beiden Formen zusammengesetzt. Aus denselben Formen ist aber die Form $(A,B,C)$ zusammengesetzt, wenn man den größten gemeinsamen Teiler von $a$, $a'$, $b+b'$ mit $=\mu$ bezeichnet und $A=\frac{aa'}{\mu\mu}$, $B\equiv b$ und $b'$ jeweils modulo $\frac{a}{\mu}$, $\frac{a'}{\mu}$. setzt,
\SourceLeaf{272}
$AC=BB-D$; somit ist diese Form der Form $F$ eigentlich äquivalent. Durch die Form $Axx+2Bxy+Cyy$ wird nun die Zahl $aa'$ dargestellt, indem man $x=\mu$, $y=0$ setzt; der größte gemeinsame Teiler dieser Werte ist $\mu$. Also kann $aa'$ auch durch die Form $F$ so dargestellt werden, dass die Werte der Unbestimmten den größten gemeinsamen Teiler $\mu$ haben (Artikel~166). Gilt also $\mu=1$, so kann $aa'$ durch die Form $F$ dargestellt werden, indem man den Unbestimmten teilerfremde Werte gibt, und diese Darstellung gehört zu dem Wert $B$ des Ausdrucks $\surd D$ (modulo $aa'$), der zu $b$, $b'$ jeweils modulo $a$, $a'$ kongruent ist. Die Bedingung $\mu=1$ gilt stets, wenn $a$, $a'$ teilerfremd sind; allgemeiner, wenn der größte gemeinsame Teiler von $a$, $a'$ zu $b+b'$ teilerfremd ist.

\Needspace{11\baselineskip}
\begin{center}
\textit{Komposition der Ordnungen.}
\end{center}

\Article{245}
\textsc{Satz.} \textit{Gehört die Form $f$ zu derselben Ordnung wie $g$ und ebenso $f'$ zu derselben Ordnung wie $g'$, so hat die Form $F$, die aus $f$, $f'$ zusammengesetzt ist, dieselbe Determinante und gehört zu derselben Ordnung wie die Form $G$, die aus $g$, $g'$ zusammengesetzt ist.}

\textit{Beweis.} Die Formen seien jeweils $f,f',F=(a,b,c)$, $(a',b',c')$, $(A,B,C)$ und ihre Determinanten $=d,d',D$. Ferner sei der größte gemeinsame Teiler von $a,2b,c$ $=m$, der größte gemeinsame Teiler von $a,b,c$ $=\mathfrak{m}$; $m'$, $\mathfrak{m}'$ sollen eine entsprechende Bedeutung für die Form $f'$ haben und $M$, $\mathfrak{M}$ für die Form $F$. Dann wird die Ordnung der Form $f$ durch die Zahlen $d,m,\mathfrak{m}$ bestimmt; daher gelten dieselben Zahlen auch für die Form $g$. Ebenso sind die Zahlen $d',m',\mathfrak{m}'$ für die Form $g'$ dieselben wie für die Form $f'$. Nach Artikel~235 werden nun die Zahlen $D,M,\mathfrak{M}$ durch $d,d',m,m',\mathfrak{m},\mathfrak{m}'$ bestimmt: $D$ ist nämlich der größte gemeinsame Teiler von $dm'm'$, $d'mm$; $M=mm'$; und $\mathfrak{M}=\mathfrak{m}\mathfrak{m}'$ (wenn zugleich $m=\mathfrak{m}$, $m'=\mathfrak{m}'$) oder $=2\mathfrak{m}\mathfrak{m}'$ (wenn $m=2\mathfrak{m}$ oder $m'=2\mathfrak{m}'$). Diese Eigenschaften von $D,M,\mathfrak{M}$ folgen daraus, dass $F$ aus $f,f'$ zusammengesetzt ist; daher erkennt man leicht, dass $D,M$ und $\mathfrak{M}$ auch für die Form $G$ gelten und somit $G$ zu derselben Ordnung wie $F$ gehört. \textit{Q.~E.~D.}

Aus diesem Grund nennen wir die Ordnung, zu der die Form $F$ gehört, zusammengesetzt aus den Ordnungen, zu denen die Formen $f,f'$ gehören. So entsteht \textit{zum Beispiel}\ aus zwei eigentlich primitiven Ordnungen stets eine ebensolche Ordnung; aus einer eigentlich primitiven und einer uneigentlich primitiven eine uneigentlich primitive. --- Ebenso ist es zu verstehen, wenn eine Ordnung aus mehreren anderen Ordnungen zusammengesetzt genannt wird.
